\documentclass[11pt]{amsart}
\usepackage[localtheorems]{raeez-math-template}

\providecommand{\C}{\mathbb{C}}
\providecommand{\Z}{\mathbb{Z}}
\providecommand{\KK}{\mathbb{K}}
\providecommand{\dbar}{\bar{\partial}}
\providecommand{\gfrak}{\mathfrak{g}}
\providecommand{\Lfrak}{\mathfrak{L}}
\providecommand{\tr}{\mathrm{tr}}
\providecommand{\gh}{\mathrm{gh}}
\providecommand{\BV}{\mathrm{BV}}
\providecommand{\KT}{\mathrm{KT}}
\providecommand{\CE}{\mathrm{CE}}
\providecommand{\chCE}{\mathrm{CE}^{\mathrm{ch}}}
\providecommand{\hCS}{\mathrm{hCS}}
\providecommand{\cA}{\mathcal{A}}
\providecommand{\cE}{\mathcal{E}}
\providecommand{\cL}{\mathcal{L}}
\providecommand{\cO}{\mathcal{O}}
\providecommand{\Obs}{\mathrm{Obs}}
\providecommand{\Sym}{\mathrm{Sym}}
\providecommand{\HH}{\mathrm{HH}}
\providecommand{\Ran}{\mathrm{Ran}}
\providecommand{\Uch}{U^{\mathrm{ch}}}
\providecommand{\barB}{\mathsf{B}}
\providecommand{\Omb}{\Omega_{\mathrm{b}}}
\providecommand{\MC}{\mathrm{MC}}
\providecommand{\kch}{\kappa_{\mathrm{ch}}}
\providecommand{\kcat}{\kappa_{\mathrm{cat}}}
\providecommand{\kBKM}{\kappa_{\mathrm{BKM}}}
\providecommand{\kfib}{\kappa_{\mathrm{fiber}}}
\providecommand{\CoHA}{\mathrm{CoHA}}
\providecommand{\Woneinf}{\mathcal{W}_{1+\infty}}
\providecommand{\Zchder}{Z^{\mathrm{der}}_{\mathrm{ch}}}

\theoremstyle{plain}
\newtheorem{theorem}{Theorem}[section]
\newtheorem{proposition}[theorem]{Proposition}
\newtheorem{lemma}[theorem]{Lemma}
\newtheorem{corollary}[theorem]{Corollary}
\theoremstyle{definition}
\newtheorem{definition}[theorem]{Definition}
\newtheorem{construction}[theorem]{Construction}
\theoremstyle{remark}
\newtheorem{remark}[theorem]{Remark}

\title{From BV-BRST to Chiral Chevalley-Eilenberg}
\author{Raeez Lorgat}
\date{2026-04-22}

\begin{document}

\begin{abstract}
Classical BV-BRST observables are Chevalley-Eilenberg cochains of
the derived gauge Lie algebra only after the Koszul-Tate resolution of
the Euler-Lagrange shell has been inserted.  On a smooth complex curve,
the Beilinson-Drinfeld chiral envelope of a perfect
$L_\infty$-algebra has its completed factorisation bar complex
identified with chiral Chevalley-Eilenberg chains.  These two facts give
comparison maps between distinct ambient categories.  The list includes
ordinary Chevalley-Eilenberg cochains, Beilinson-Drinfeld chiral
chains, factorisation homology, derived centres, Drinfeld doubles,
boundary current algebras, and BV observables.  For six-dimensional
holomorphic Chern-Simons on $\C^3$, the local classical dg Lie algebra
is $\gfrak\otimes\Omega^{0,\bullet}_c(\C^3)$; after choosing a
projection to a curve and a transverse residue model, its residue
associated graded is an $E_1$-chiral CE object on the curve.  The
quantum comparison requires the Costello effective master equation,
$\hbar$-adic completion, quartic anomaly trivialisation, factorisation
descent, and the chosen propagator.  For compact CY$_3$ inputs, the
corresponding statement requires an orientation and determinant-line
framing, elliptic gauge fixing, a regulator, analytic completion, TCFT
or factorisation gluing, the QME, and the filtered
$\HH^{-2}_{E_1}(A,A)$ obstruction comparison.  The specialised CY$_3$
output algebra is $E_1$-chiral; the $E_2$ structure belongs to the
centre, double, or module layer under its own hypotheses.
\end{abstract}

\maketitle
\tableofcontents

\section{Classical CE and the Koszul-Tate Shell}
\label{sec:linear}

Let $\KK$ be a field of characteristic zero.  A finite-dimensional Lie
algebra $\Lfrak$ has Chevalley-Eilenberg chains
\[
  \CE_*(\Lfrak)
  =
  \bigl(\Sym(\Lfrak[1]), d_\CE\bigr),
  \qquad
  d_\CE(x\wedge y) = [x,y],
\]
where $d_\CE$ is extended as a coderivation of the cocommutative
coalgebra.  For a module $M$,
\[
  \CE^*(\Lfrak;M)
  =
  \bigl(\Sym(\Lfrak^*[-1])\otimes M,d_\CE\bigr),
\]
with the usual alternating sum of the Lie bracket and the action on
$M$.  The identity $d_\CE^2=0$ is Jacobi.

Let $F$ be the space of classical fields, let $\gfrak$ be a closed
irreducible gauge Lie algebra acting by infinitesimal gauge
transformations, and let $S_0\in\cO(F)^{\gfrak}$ be the bare action.
The BV graded manifold is
\[
  \cE^\BV =
  T^*[-1](F/\!/\gfrak)
  =
  \gfrak[1]\oplus F\oplus F^*[-1]\oplus\gfrak^*[-2].
\]
The four summands are ghosts, fields, antifields, and antighosts, with
ghost numbers $1,0,-1,-2$.  The degree $-1$ symplectic form is the
natural contraction
$F\otimes F^*[-1]\to\KK[-1]$ together with
$\gfrak[1]\otimes\gfrak^*[-2]\to\KK[-1]$.  The classical BV action is
\[
  S_{\mathrm{cl}}
  =
  S_0[\phi]
  + \phi^\dagger(X_c\phi)
  + \frac{1}{2}c^\dagger[c,c],
\]
and the classical master equation is equivalent to gauge invariance of
$S_0$ on the shell and Jacobi for $\gfrak$.

Let $\Sigma=\{dS_0=0\}\subset F$ and
$I_\Sigma=(\partial_i S_0)$.  The Koszul-Tate complex
\[
  \KT_* =
  \Sym_{\cO(F)}(F^*[-1]),
  \qquad
  d_\KT\phi_i^\dagger=\partial_iS_0,
\]
resolves $\cO(\Sigma)$ when the equations of motion form a regular
sequence.  If the shell is derived or non-regular, $\KT_*$ is replaced
by a cofibrant Koszul-Tate resolution of the derived critical locus.
The gauge action extends to $\KT_*$, so one obtains the
BRST-Koszul-Tate bicomplex
\[
  \CE^*(\gfrak;\KT_*)
  =
  \Sym(\gfrak^*[-1])
  \otimes
  \Sym(F^*[-1])
  \otimes
  \cO(F),
  \qquad d=d_\CE+d_\KT .
\]

\begin{theorem}[BV observables and the Koszul-Tate CE complex]
\label{thm:bv-ce-ktshell}
Assume $\Sigma=\{dS_0=0\}$ is a complete intersection, the strict
$\gfrak$-action preserves $S_0$, and the displayed $\KT_*$ is the chosen
Koszul-Tate resolution of the shell.  Then
\[
  \bigl(\cO(\cE^\BV),\{S_{\mathrm{cl}},-\}\bigr)
  \simeq
  \bigl(\CE^*(\gfrak;\KT_*),d_\CE+d_\KT\bigr)
\]
as differential graded commutative algebras.  In ghost number zero the
cohomology is $\cO(\Sigma)^{\gfrak}$.  For reducible or open gauge
theories the same statement is obtained by replacing $\gfrak$ with the
derived gauge $L_\infty$-algebra and $\KT_*$ with the full
Koszul-Tate resolution, including higher Tate generators for reducible
symmetries.
\end{theorem}

\begin{proof}
The Hamiltonian vector field of $S_{\mathrm{cl}}$ decomposes as
\[
  \{S_0,-\}
  +\{c^\alpha X_\alpha,-\}
  +\left\{\frac{1}{2}c^\dagger[c,c],-\right\}.
\]
The first summand is $d_\KT$, the second is the CE differential with
coefficients in $\KT_*$, and the third is the CE differential of
$\gfrak$.  The anticommutation relation is gauge invariance of $S_0$.
The square-zero conditions are regularity of the shell and Jacobi.
This is the standard BV-BFV comparison in the form used by
Henneaux-Teitelboim, Stasheff, and Costello-Gwilliam.
\end{proof}

\begin{remark}[Derived gauge Lie algebra]
\label{rmk:derivedlie}
For a closed strict gauge algebra the data may be represented by a dg
Lie algebra.  For open gauge algebras and field theories whose gauge
commutators close only modulo the equations of motion, the natural
object is an $L_\infty$-algebra.  The CE algebra of that derived gauge
algebra presents the derived critical quotient
$\mathbb{R}\mathrm{Crit}(S_0)/\gfrak$ only after the Koszul-Tate shell
is part of the coefficient data.  Without the Koszul-Tate generators
one computes off-shell gauge cohomology.  Holomorphic Chern-Simons is
strict at the classical local level: its derived Lie algebra is
$\Omega^{0,\bullet}(-,\gfrak)$ with $\dbar$ and the pointwise bracket.
\end{remark}

\section{Chiral Envelope and Chiral CE}
\label{sec:chiral}

Let $C$ be a smooth complex curve and let $\cL$ be a perfect,
bounded-below sheaf of $\cD_C$-linear $L_\infty$-algebras on $C$, or a
pro-nilpotent filtered completion of such a sheaf.  The model
$\gfrak\otimes\Omega^{0,\bullet}_c(C)$ with $\dbar$-differential is the
Dolbeault realisation.  The Beilinson-Drinfeld chiral envelope
$\Uch(\cL)$ is an augmented factorisation $\cD$-module on $\Ran(C)$ with
chiral operation $\mu_{\mathrm{ch}}$.  On a disk $D$ with coordinate
$z$, its local sections recover the enveloping algebra of the
corresponding loop $L_\infty$-algebra subject to the BD chiral
relations:
\[
  \Uch(\cL)(D)
  \simeq
  U\bigl(\cL(D)\otimes\KK(\!(z)\!)\bigr)/(\text{chiral relations}).
\]
The disk calculation records the associative shadow.  The Ran-space
factorisation structure records the operator product at colliding
points.

\begin{definition}[Chiral CE chains]
\label{def:chce}
The chiral Chevalley-Eilenberg chains of $\cL$ are the factorisation
$\cD$-module whose fibre over a finite configuration
$\{x_1,\ldots,x_n\}\subset C$ is the symmetric coalgebra on
$\cL[1]$ at those points, with differential equal to the pointwise CE
differential plus the chiral fusion terms.  On a disk,
\[
  \chCE_*(\cL)(D)
  =
  \bigl(\Sym(\cL(D)[1]),d_\CE+d_{\mathrm{fus}}\bigr),
\]
where $d_{\mathrm{fus}}$ is induced by OPE residues.
\end{definition}

\begin{theorem}[Factorisation bar and chiral CE]
\label{thm:chce-bar}
For a perfect bounded-below $\cD$-linear $L_\infty$-algebra on a smooth
curve, with augmented chiral envelope and conilpotent completed
factorisation bar construction, there is a natural quasi-isomorphism of
factorisation $\cD$-modules
\[
  \widehat{\barB}(\Uch(\cL))
  \simeq
  \widehat{\chCE}_*(\cL).
\]
On a disk it reduces to Quillen's quasi-isomorphism
$\barB(U(\Lfrak))\simeq\CE_*(\Lfrak)$ applied to the local loop
$L_\infty$-algebra.
\end{theorem}

\begin{proof}
The chiral envelope is defined by the universal chiral operations of
$\cL$.  Its factorisation bar complex is the cosimplicial object whose
face maps insert those operations along collisions in $\Ran(C)$.  These
face maps are exactly the fusion summands in the chiral CE differential.
The local disk reduction is Quillen's computation of the bar complex of
an enveloping algebra.  Ran compatibility is the content of
Beilinson-Drinfeld chiral homology and of the
Francis-Gaitsgory chiral Koszul duality theorem; the
$L_\infty$ form is the Costello-Gwilliam factorisation version.
\end{proof}

\begin{remark}[Completion]
\label{rmk:bar-ce-completion}
For Dolbeault and current-algebra models the bar and CE sides are the
same filtered completion.  Without the augmentation, conilpotence, and
completion hypotheses, the displayed comparison is a disk-level
calculation.  A functorial equivalence of factorisation \(\cD\)-modules
on \(\Ran(C)\) requires those hypotheses.
\end{remark}

\begin{remark}[$E_1$ and $E_2$ scope]
\label{rmk:E1E2scope}
On a curve the chiral algebra is $E_1$-ordered along the complex
direction.  Complex surfaces can carry genuine $E_2$-chiral objects.
For a CY$_3$ specialisation to a curve, the output algebra $A$ is
$E_1$; the $E_2$-braiding is recovered, when the hypotheses are met, on
$\mathcal Z(\mathrm{Rep}^{E_1}(A))$ or on the corresponding Drinfeld
double.  Thus \(A\) carries the specialised \(E_1\)-chiral product, and
the braided operation belongs to the centre or double layer.
\end{remark}

\section{Master Equations and CE Cocycles}
\label{sec:master}

Let $\cL$ be an $L_\infty$-algebra and let $\Omb$ be a unital
commutative dg algebra.  A Maurer-Cartan element
$\alpha\in\cL\otimes\Omb$ satisfies
\[
  \MC(\alpha)
  =
  \sum_{k\geq 1}\frac{1}{k!}
  \ell_k(\alpha^{\otimes k})
  =
  0.
\]
Equivalently, the CE differential on $\CE^*(\cL;\Omb)$ annihilates the
corresponding exponential class.  In the classical BV setting this is
the classical master equation for the derived gauge algebra of
Remark~\ref{rmk:derivedlie}.

Costello's quantum BV formalism replaces the classical differential by
the scale-dependent effective equation
\[
  QI[L]
  +\frac{1}{2}\{I[L],I[L]\}_L
  +\hbar\,\Delta_L I[L]
  =
  0,
\]
where $Q$ is the linear BRST operator, $\{-,-\}_L$ and $\Delta_L$ are
defined by the scale-$L$ propagator, and $I[L]$ is the effective
interaction.  This equation is meaningful only in the renormalised,
$\hbar$-adically completed algebra of local functionals.

\begin{proposition}[Quantum transfer to CE form]
\label{prop:qme-hce}
Let $(\cE,Q,\omega)$ be a perturbative BV theory with a chosen
propagator $P_{\varepsilon,L}$ and an anomaly-free effective family
$I[L]$ solving the quantum master equation in the completed local
functional algebra.  A filtered deformation retract of $(\cE,Q)$ onto
its cohomology transfers the effective interaction to operations
$\ell^\hbar_k$ on the cohomology, given by connected Feynman graphs
with propagator edges and powers of $\hbar$ determined by loop number.
The transferred interaction is a Maurer-Cartan element for the
completed quantum CE differential determined by $\{\ell^\hbar_k\}$.
\end{proposition}

\begin{proof}
The homological perturbation lemma applies to the contraction data and
expands the transferred differential as the sum over connected graphs.
Tree graphs give the classical transferred $L_\infty$ operations; loop
graphs contribute the powers of $\hbar$ and the BV Laplacian insertions.
Costello's effective master equation is precisely the compatibility
condition ensuring that the transferred differential squares to zero.
\end{proof}

\begin{remark}[Chiral quantum comparison]
\label{rmk:chiral-qme}
On a curve, a chiral quantum comparison is a map
\[
  \Theta_C:
  \Obs^q_{\BV}(C)
  \longrightarrow
  \widehat{\chCE}_*(\cL)[[\hbar]]
\]
to the completed chiral CE complex.  Its associated graded at
$\hbar=0$ is the classical factorisation-bar comparison of
Theorem~\ref{thm:chce-bar}.  The quantum map depends on the chiral
Green kernel, the renormalisation scheme, and the anomaly cancellation
data.  A binary chiral product equation has too little structure:
the BV Laplacian, the propagator, and the local obstruction class have
disappeared.
\end{remark}

\begin{proposition}[Quartic obstruction and two-point counterterm]
\label{prop:quartic-anomaly-counterterm}
Let $\rho:\gfrak\to\operatorname{End}(V)$ be a representation.  Define
\[
  I_4^\rho(x_1,x_2,x_3,x_4)
  =
  \frac{1}{4!}\sum_{\sigma\in S_4}
  \operatorname{Tr}_V\!\left(
    \rho(x_{\sigma(1)})\rho(x_{\sigma(2)})
    \rho(x_{\sigma(3)})\rho(x_{\sigma(4)})
  \right),
  \qquad
  I_4^\rho\in\Sym^4(\gfrak^\vee)^\gfrak .
\]
In complex dimension three the one-loop QME obstruction for nonabelian
hCS is represented, up to the standard propagator normalisation, in the
Costello local Chevalley complex by
\[
  \Theta^\rho_U(\alpha_0,\alpha_1,\alpha_2,\alpha_3)
  =
  \int_U
  \left[
  I_4^\rho\bigl(
    \alpha_0,\partial\alpha_1,\partial\alpha_2,\partial\alpha_3
  \bigr)
  \right]_{(3,3)},
  \qquad
  [\hbar\Theta^\rho]\in H^1_{\mathrm{loc}}.
\]
For a basis $\{T_a\}$ of $\gfrak$ and dual basis $\{T^a\}$ for the
chosen invariant pairing, the quadratic Casimir
\[
  C_2^\rho=\sum_a\rho(T_a)\rho(T^a)
\]
is the colour factor of the two-point kinetic counterterm.  The quartic
class has degree \(1\) and arity \(4\); the quadratic counterterm has
degree \(0\) and arity \(2\).  The latter cannot trivialise the former.
\end{proposition}

\begin{proof}
The Costello-Li one-loop obstruction graph in complex dimension three
is the four-wheel.  Its Lie factor is the symmetrised trace of four
inputs, and the form factor supplies three holomorphic derivatives,
which produce the displayed \((3,3)\)-density.  Costello's obstruction
theory places this local functional in degree one, the first quantum
slot where the scale-\(L\) master equation can fail.  Closing two legs
of the cubic hCS vertex gives instead the contraction
\(\sum_a\rho(T_a)\rho(T^a)\); the remaining tensor type is the kinetic
functional \(\int\Omega\wedge\operatorname{Tr}_\rho(\cA\,\dbar\cA)\).
This is a degree-zero field-strength renormalisation.
\end{proof}

\section{Six-Dimensional Holomorphic Chern-Simons}
\label{sec:reduction}

Let $X=\C^3$ with holomorphic volume form
$\Omega=dz_1\wedge dz_2\wedge dz_3$.  Six-dimensional holomorphic
Chern-Simons with structure Lie algebra $\gfrak$ and invariant
pairing $\tr$ has BV fields
\[
  \cE^\BV_{\hCS}
  =
  \Omega^{0,\bullet}(X,\gfrak)[1],
\]
where the invariant pairing identifies the antifields with the
Dolbeault degrees \(2\) and \(3\).  Without this identification the
dual summand is \(\Omega^{3,\bullet}_c(X,\gfrak^\vee)[-2]\), paired
with fields by \(\Omega\).
Unwound by Dolbeault degree, the fields are
$c\in\Omega^{0,0}$, $\cA\in\Omega^{0,1}$,
$\cA^\dagger\in\Omega^{0,2}$, and
$c^\dagger\in\Omega^{0,3}$.  The degree $-1$ pairing is
\[
  \omega(\alpha,\beta)
  =
  \int_X\Omega\wedge\tr(\alpha\wedge\beta).
\]
The classical action is
\[
  S_{\mathrm{cl}}
  =
  \int_X\Omega\wedge\tr\!\left(
  \frac{1}{2}\cA\dbar\cA
  +\frac{1}{6}\cA[\cA,\cA]
  +\cA^\dagger\dbar c
  +\cA^\dagger[\cA,c]
  +\frac{1}{2}c^\dagger[c,c]
  \right).
\]
The underlying classical dg Lie algebra is
\[
  \Lfrak_X
  =
  \gfrak\otimes\Omega^{0,\bullet}_c(X),
  \qquad
  d=\dbar,
  \qquad
  \ell_2=[-,-]_{\gfrak}\otimes\wedge ,
\]
with no higher classical brackets.  Theorem~\ref{thm:bv-ce-ktshell}
specialises to
\[
  \bigl(\cO_{\mathrm{loc}}(\cE^\BV_{\hCS}),\{S_{\mathrm{cl}},-\}\bigr)
  \simeq
  \widehat{\CE}^*(\Lfrak_X),
\]
the completed continuous CE cochains of the compactly supported local
dg Lie algebra.

\subsection{Transverse Residue Model}

Choose the projection
$\pi:\C^3\to C=\C$, $(z_1,z_2,z_3)\mapsto z_1$.  The raw derived
push-forward \(R\pi_!\Omega^{0,\bullet}_c(\C^3)\) retains the
transverse Dolbeault complex of \(\C^2\).  The Dolbeault complex of
\(C\) appears after
choosing the local-cohomology summand supported on the transverse
origin and the Bochner-Martinelli residue functional
\[
  \operatorname{Res}_{\C^2,0}:
  \Omega^{0,\bullet}_{c,\{0\}}(\C^2)
  \longrightarrow
  \C[-2].
\]
Applying \(\operatorname{Res}_{\C^2,0}\) and the conventional CY shift
gives the residue curve dg Lie algebra
\[
  \cL_C^{\mathrm{res}}
  =
  \gfrak\otimes\Omega^{0,\bullet}_c(C).
\]
The residue current, support condition, filtration, and completion are
part of the datum.

\begin{proposition}[Local classical comparison on \texorpdfstring{$\C^3$}{C3}]
\label{prop:c3-classical-comparison}
For $X=\C^3$ with the projection above and the transverse residue model
fixed, the residue associated graded of classical hCS observables maps
by a quasi-isomorphism to the chiral CE complex
\[
  \mathrm{gr}_{\mathrm{res}}\,\pi_*\Obs_{\hCS}^{\mathrm{cl}}
  \simeq
  \chCE_*\bigl(\cL_C^{\mathrm{res}}\bigr).
\]
The resulting chiral algebra on $C$ is $E_1$.
\end{proposition}

\begin{proof}
The BV side is the CE algebra of $\Lfrak_X$.  The relative compact
support calculation is followed by the residue projection
\(\operatorname{Res}_{\C^2,0}\), which passes to the associated graded
for the transverse-support filtration and produces
\(\cL_C^{\mathrm{res}}\).  The
Beilinson-Drinfeld chiral envelope and Theorem~\ref{thm:chce-bar}
then identify the factorisation bar with
$\chCE_*(\cL_C^{\mathrm{res}})$.  Since the final factorisation
direction is the curve direction, the native chiral order is $E_1$.
\end{proof}

\begin{proposition}[Quantum comparison map]
\label{prop:c3-quantum-comparison}
Assume the Costello-Williams heat-kernel propagator on $\C^3$, the
residue filtration above, factorisation descent for the projection
\(\pi\), an effective family solving the QME after the quartic anomaly
class of
Proposition~\ref{prop:quartic-anomaly-counterterm} has been
trivialised, and the $\hbar$-adic completion of the local functional
algebra.  Then there is a filtered comparison map
\[
  \Theta_\pi:
  \widehat{\pi_*\Obs^q_{\hCS}}^{\,\mathrm{res}}
  \longrightarrow
  \widehat{\chCE}_*(\cL_C^{\mathrm{res}})[[\hbar]]
\]
whose associated graded at $\hbar=0$ is the classical comparison of
Proposition~\ref{prop:c3-classical-comparison}.  A quantum
quasi-isomorphism is an additional theorem requiring the stated
renormalisation and anomaly hypotheses.
\end{proposition}

\begin{proof}
The scale-dependent propagator defines the BV bracket and Laplacian on
the completed local functional algebra.  Proposition~\ref{prop:qme-hce}
transfers the effective interaction to graph-weighted operations on the
cohomology.  The chiral Green kernel on the curve gives the target
completed chiral CE differential.  The associated graded statement is
the classical comparison.  The filtered statement follows only after
the anomaly term vanishes and the completions agree.
\end{proof}

\begin{remark}[Compact CY$_3$ scope]
\label{rmk:compact-cy3-scope}
For $X=K3\times E$ or a non-formal compact CY$_3$, the projection to a
curve retains the full Dolbeault and Mukai fibre data.  The specialised
output is an $E_1$-chiral algebra on the curve only after the
specialisation cycle and analytic data are fixed.  The $E_2$ structure
belongs to the Drinfeld centre or double of its representation
category.  The compact hCS statement requires an orientation and
determinant-line framing for the BV pairing, elliptic gauge fixing, a
regulator, completed local functionals, a QME solution, trivialisation
of the quartic \(I_4^\rho\) anomaly by local or closed-sector data,
TCFT or factorisation sewing, and the filtered
$\HH^{-2}_{E_1}(A,A)$ obstruction comparison.  These are hypotheses,
not formal consequences of the equation \(c_1(X)=0\).
\end{remark}

\section{Separated Objects}
\label{sec:separated}

The dictionary uses distinct constructions.
\begin{center}
\begin{tabular}{p{0.30\linewidth}p{0.62\linewidth}}
\(\CE^*(\Lfrak)\) &
ordinary dg commutative algebra of cochains on a derived gauge
Lie algebra; for BV observables the Koszul-Tate shell is part of the
coefficient data \\
\(\chCE_*(\cL)\) &
factorisation \(\cD\)-module on \(\Ran(C)\), with fusion differential \\
\(\int_C \Uch(\cL)\) &
factorisation homology or chiral homology, a global object obtained
after descent and support choices \\
\(\Obs^{\mathrm{cl}}_{\BV}\) &
classical BV factorisation algebra of local observables on the derived
Euler-Lagrange shell \\
\(\Obs^q_{\BV}\) &
\(\hbar\)-complete quantum BV factorisation algebra after the QME \\
\(\Zchder(A)\) &
derived chiral centre or Hochschild-type bulk object of an \(E_1\)
chiral algebra \(A\) \\
\(D(Y^+)\) &
Drinfeld double pairing positive and negative Hall halves \\
\(\Uch(\gfrak((z)))\) &
boundary current algebra, with central extension fixed by loop data
\end{tabular}
\end{center}
The maps among them are comparison maps with hypotheses.  Each map
names its source, target, topology, completion, and gluing datum.

\begin{proposition}[Antifields and the shell]
\label{prop:antifields-shell}
Removing antifields replaces
$\CE^*(\gfrak;\KT_*)$ by $\CE^*(\gfrak;\cO(F))$.  The latter computes
gauge cohomology on all fields, whereas the former computes gauge
cohomology on the derived shell.  The antifields are exactly the
Koszul-Tate generators cutting out $\Sigma=\{dS_0=0\}$.
\end{proposition}

\begin{proposition}[Strict and derived gauge algebras]
\label{prop:strict-derived}
Holomorphic Chern-Simons has a strict classical dg Lie algebra
$\Omega^{0,\bullet}(-,\gfrak)$ with $\dbar$ and the wedge bracket.
Open gauge theories require $L_\infty$ brackets.  Quantum effective
operations obtained from graph transfer are operations on the
renormalised effective theory; the strict classical hCS dg Lie algebra
keeps only \(\dbar\) and the wedge-Lie bracket.
\end{proposition}

\begin{proposition}[Bar, duals, centre, fields, and polyvectors]
\label{prop:five-objects-separation}
For an \(E_1\)-chiral algebra \(A\), the five symbols
\[
  A,\qquad B(A),\qquad A^i,\qquad A^!,\qquad \Zchder(A)
\]
have different definitions.  The identity \(\Omega B(A)\simeq A\)
expresses bar-cobar inversion.  The object \(A^i\) is the Koszul dual
coalgebra.  The object \(A^!\) is the Verdier or Koszul dual algebra in
the chosen duality context.  The object \(\Zchder(A)\) is the bulk or
braided layer at \(d\geq3\).

The hCS field complex
\(\Omega^{0,\bullet}(X,\operatorname{ad}P)[1]\), its BV antifields
paired by \(\int_X\Omega_X\wedge\operatorname{Tr}(-\, -)\), and the
Hochschild polyvectors of a CY category are separate objects.  For
example,
\[
  \HH^\bullet(\operatorname{Perf}(\C^3))
  \simeq
  \C[x,y,z]\otimes
  \bigwedge\nolimits^\bullet
  \langle\partial_x,\partial_y,\partial_z\rangle .
\]
These polyvectors control categorical deformations and centre
constructions.  A finite-dimensional hCS gauge algebra enters as an
external coefficient or framing datum.
\end{proposition}

\begin{proposition}[Bulk dependence at loop level]
\label{prop:bulk-loop}
The classical curve model remembers only the transverse residue class.
The quantum comparison remembers the six-dimensional bulk through the
heat-kernel propagator, the invariant pairing, and the obstruction
cocycle.  Boundary current-algebra central extensions are therefore
loop-level data, beyond the ordinary CE complex alone.
\end{proposition}

\section{Consistency Constraints}
\label{sec:cross}

\begin{itemize}
\item \textbf{Class M bar rank.} The assertion that
$\barB(\Uch(\mathrm{Vir}))$ has cohomology rank $6^g$ is a
cohomology-level computation.  The factorisation-bar-to-chiral-CE
mechanism supplies the comparison surface; the per-genus six-generator
calculation is a separate input.

\item \textbf{CoHA of $\C^3$.} The cohomological Hall algebra of
$\C^3$ is $Y^+(\widehat{\mathfrak{gl}}_1)$, the positive half.  The full
\(\Woneinf\) or full affine Yangian requires the double, centre,
pairing, completion, and evaluation constructions:
\[
  \CoHA(\C^3)\cong Y^+(\widehat{\mathfrak{gl}}_1),
  \qquad
  \Woneinf\ \text{appears only after}\quad
  Y^+\mapsto D(Y^+)\mapsto\text{centre/evaluation}.
\]

\item \textbf{Product Euler characteristic.}
$\kcat(K3\times E)=\chi(\cO_{K3})\chi(\cO_E)=2\cdot 0=0$.  The K3 fibre
value $\chi(\cO_{K3})=2$ is a fibre witness; the total-space value is
zero.

\item \textbf{K3 \(\times\) E invariants.} For \(X=K3\times E\),
\[
  \kcat(X)=0,\qquad
  \kch^{\mathrm{compact}}(X)=0,\qquad
  \kch^{\mathrm{Heis}}(X)=3,\qquad
  \kBKM(\mathfrak g_{\Delta_5})=5,\qquad
  \kfib(X)=24.
\]
The sources are respectively the total-space Euler characteristic, the
compact Hodge supertrace, the Heisenberg-Mukai specialisation,
Gritsenko's \(\Delta_5\) Borcherds denominator with \(c_1(0)/2=5\), and
the K3 Mukai lattice rank.  The proposed expression
\(\kBKM=\kch^{\mathrm{compact}}+\chi(\cO_E)\) has right-hand side
zero, while the Borcherds weight is five.

\item \textbf{Compact Hodge supertrace.} On compact CY$_d$ inputs where
the comparison is proved,
\[
  \kch(A_X)=\sum_q(-1)^q h^{0,q}(X)=\chi(\cO_X)=\kcat(X).
\]
This is a theorem on the compact constructed locus.  Singular and
non-compact examples require their own construction.

\item \textbf{CY$_3$ operadic level.} For $d\geq 3$ the specialised
output on the curve is an \(E_1\)-chiral algebra \(A\).  The braided
structure is a centre or double structure on
\(\mathcal Z(\mathrm{Rep}^{E_1}(A))\) under its own hypotheses.

\item \textbf{Hall-Borcherds and self-mirror targets.} The
\(K3\times E\) Hall-Borcherds double and the Mukai self-mirror K3
Yangian branch are separate targets until an explicit centre or double
comparison functor has been specified.
\end{itemize}

\section{Open Problems}
\label{sec:open}

\begin{enumerate}[label=\textup{(\arabic*)}]
\item Construct the pushed-forward $L_\infty$ sheaf for
$K3\times E\to E$ from the full Dolbeault and Mukai fibre data, and
compare it with the Hall-Borcherds double or the K3 Yangian branch with
the target fixed in advance.
\item Prove the quantum chiral comparison map at genus $g\geq 1$ with
the BV Laplacian, curvature term, and Positselski coderived completion
all present.
\item Upgrade Proposition~\ref{prop:c3-quantum-comparison} to a
homotopy-coherent comparison of factorisation algebras from the
completed disk-complex comparison.
\item Formulate the higher-dimensional Ran analogue of
Theorem~\ref{thm:chce-bar} for surfaces and higher bases separately
from the curve case.
\end{enumerate}

\end{document}
